\section{Methods (Search and Selection)}

\noindent
\textbf{Protocol adherence.}  
All screening, selection, extraction, and second-reader verification procedures
followed a predefined protocol (Version~0.2), finalised prior to data
collection. The protocol specifies the source universe, eligibility criteria,
search strategy, DocID conventions, evidence-capture rules, and second-reader
sampling plan. The full protocol is reproduced in Appendix~B and archived with
the repository artefacts.


\subsection{Sources}

Two categories of sources informed the evidence base for this review:

\begin{itemize}
    \item \textbf{Fixed UK sources:} Office for Students (OfS); Quality Assurance Agency for Higher Education (QAA); Department for Education (DfE); UK~ENIC (ECCTIS); Universities UK (UUK); SCQF Partnership; Scottish Funding Council (SFC); Welsh Government / CTER; Department for the Economy Northern Ireland (DfE~NI); Ofqual.
    \item \textbf{Context sources (Sri Lanka):} Sri Lanka Qualifications Framework 
    (SLQF) / University Grants Commission (UGC); British Council Sri Lanka.
\end{itemize}

All sources correspond to entries in \texttt{screening\_log\_\_v0.1\_freeze\_2025-11-24.csv} (12 items).

\subsection{Date Window}

The search prioritised current versions in force as of 2025, with backwards lineage consulted to approximately 2008 where required to understand regulatory transitions or framework evolution.

\subsection{Search Queries}

Searches used eight predefined query groups executed verbatim on Google and source-domain search bars:

\begin{itemize}
    \item G1) \texttt|site:qaa.ac.uk ("FHEQ" OR "qualification descriptors" OR "levels")|
    \item G2) \texttt|site:officeforstudents.org.uk (regulatory OR "conditions of registration" OR "validation" OR "franchise" OR "transnational education")|
    \item G3) \texttt|site:ecctis.com (recognition OR comparability OR "GCE A/L" OR "Sri Lanka")|
    \item G4) \texttt|site:gov.uk ("Quality Code" OR "sector-agreed principles" OR "higher education regulation")|
    \item G5) \texttt|site:scqf.org.uk (SCQF OR "level descriptors" OR "credit")|
    \item G6) \texttt|site:universitiesuk.ac.uk ("transnational education" OR TNE OR "quality assurance")|
    \item G7) \texttt|("exit award" OR "intermediate award" OR CertHE OR DipHE OR "graduate diploma" OR "postgraduate diploma") site:.ac.uk filetype:pdf|
    \item G8) \texttt|("validation" OR "franchise") ("partner" OR "delivery") ("quality assurance" OR "oversight") site:.ac.uk|
\end{itemize}

\subsection{Inclusion Criteria}

In line with Protocol~v0.2, items were included where they met one or more of the following:

\begin{itemize}
    \item Official frameworks, award descriptors, and qualification statements.
    \item Regulatory frameworks, conditions, sector-recognised principles.
    \item University or awarding-body regulations for validation, franchise, or TNE.
    \item Official guidance/handbooks for partnerships and external quality assurance.
    \item Authoritative sector reports and peer-reviewed analyses relevant to TNE or recognition.
\end{itemize}

\subsection{Exclusion Criteria}

Items were excluded if they met any of the following:

\begin{itemize}
    \item Marketing, opinion, consultancy, or agency content.
    \item Non-UK or non-higher education materials.
    \item Superseded sources lacking lineage value.
    \item General consumer information without regulatory or standards relevance.
    \item Items with no identifiable year, authoritativeness, or provenance.
\end{itemize}

\subsection{Screening Process}

Screening followed a defined workflow:

\begin{itemize}
    \item Each record was logged in \texttt{screening\_log\_\_v0.1\_freeze\_2025-11-24.csv} under a DocID of the form \texttt{SRC-YY-\#\#\#}.
    \item Items were assigned a status: Include / Exclude / Borderline.
    \item ReasonCodes were applied per dictionary: I1--I4 (Inclusion), E1--E4 (Exclusion), B1 (Borderline).
    \item Where licensing permitted, PDFs or web-prints of included items were downloaded locally for analysis. For copyright reasons, the public artefact bundle includes only the metadata in \texttt{evidence\_manifest\_\_v0.1\_freeze\_2025-11-24.csv}, not the PDFs themselves.
    \item Second-reader sampling covered: random $\sim$15\% of all screened items, all Borderline items, and all Excluded items.
    \item \raggedright Disagreements were resolved by discussion; escalated cases were recorded in \texttt{agreement\_summary\_\_v0.1\_2025-11-24.csv}.
\end{itemize}

\subsection{Data Extraction}

Data extraction used the template in \texttt{evidence\_table\_\_v0.1\_freeze\_2025-11-24.csv}, with the following fields:

\begin{itemize}
    \item DocID; Title; Year; Jurisdiction; DocType; PrimaryFocus;
    \item KeyProvisions/Descriptors; Levels/Credits; ExitAwards;
    \item TNE/PartnerRules; Recognition/Entry (ENIC);
    \item Citations; UseInReview.
\end{itemize}

\subsection{Risk and Limitations}
\begin{itemize}
    \item Search-process limitations were documented in real time and include the behaviours described below.
    
    \item Risks included regulatory updates between capture and analysis, site restructuring, PDFs without pagination, and sector reports that describe practice rather than prescribe requirements.
    \item A small number of predefined queries (G3, G4 and G7) returned atypical results (e.g., image-only outputs or no indexed hits). These outcomes were retained verbatim for transparency, consistent with PRISMA expectations. As recommended in the protocol, no post hoc modification of queries was undertaken; results were recorded as returned, with additional checks made via internal domain search bars where applicable.

    
\end{itemize}


\subsection{Deviations From Protocol}

None.
